\documentclass[11pt,letterpaper]{article}

\usepackage[utf8]{inputenc}
\usepackage[T1]{fontenc}
\usepackage[spanish,es-noquoting]{babel}
\usepackage{lmodern}
\usepackage[margin=2.5cm]{geometry}

\usepackage{microtype}
\usepackage{graphicx}
\usepackage{array}
\usepackage{tabularx}
\usepackage{booktabs}
\usepackage{tikz}
\usepackage{enumitem}
\usepackage{listings}
\usepackage[hidelinks]{hyperref}

\usetikzlibrary{arrows.meta,positioning,calc}

\setlength{\parindent}{0pt}
\setlength{\parskip}{6pt}
\setlength{\emergencystretch}{2em}

\setlist{nosep,leftmargin=1.4em}

\pagestyle{plain}

\lstset{
    basicstyle=\ttfamily\small,
    breaklines=true,
    columns=fullflexible,
    keepspaces=true,
    showstringspaces=false
}

\newcommand{\archivo}[1]{\texttt{\detokenize{#1}}}

\newcommand{\crc}[3]{%
    \fbox{%
        \begin{minipage}[t]{.445\textwidth}
            \small
            \textbf{#1}

            \textbf{Responsabilidades}

            #2

            \textbf{Colaboradores}

            #3
        \end{minipage}%
    }%
}

\tikzset{
    modulo/.style={
        draw,
        align=center,
        font=\small,
        minimum height=1cm,
        text width=4cm
    },
    flecha/.style={
        -{Stealth[length=2mm]}
    },
    etiqueta/.style={
        font=\scriptsize,
        fill=white,
        inner sep=2pt
    }
}

\begin{document}

\begin{center}
    \textbf{Universidad Nacional Autónoma de México}\\
    Facultad de Ciencias

    \textbf{Reporte de Modelado y Programación}\\
    Chat con servidor en C y cliente en C\#

    Guillermo Ricardo Gómez Gómez
\end{center}

\section{Resumen}

Desarrollé un sistema de chat con un servidor en C y un cliente de
consola en C\#, utilizando TCP, cJSON, .NET y
\archivo{System.Text.Json}. El proyecto fue mi primer acercamiento
a ambos lenguajes y tuve que aprender sobre administración de memoria,
comunicación mediante sockets e interpretación de mensajes JSON.
La principal dificultad de diseño fue la acumulación de responsabilidades
en \archivo{servidor.c}. Para resolverla, separé el manejo de conexiones,
el protocolo y la ejecución de operaciones. Como resultado, obtuve una
versión funcional del chat básico con identificación, consulta de
usuarios, estados, mensajes públicos y privados, y notificaciones de
conexión y desconexión.

\section{Introducción}

El objetivo del proyecto fue desarrollar un servidor en C y un cliente
en C\# que se comunicaran mediante el protocolo definido para la materia.
Además de obtener una aplicación funcional, tuve que aprender a
organizar un programa que fue creciendo conforme incorporaba nuevas
operaciones. En este reporte explico las decisiones que tomé, las
dificultades que encontré y los cambios que necesité hacer para
continuar el desarrollo.

Como mi experiencia previa era principalmente en Java, trabajar con
C\# me resultó más fácil que trabajar con C, sobre todo por el manejo
de la memoria. En C comencé con un programa mínimo y el manejo de
argumentos para recibir el puerto del servidor. Después tuve que
comprender el uso de apuntadores, la reserva y liberación de memoria,
y la separación entre cabeceras e implementaciones. Para el cliente
necesité familiarizarme con la estructura de un proyecto de .NET y
con la coordinación de la entrada del usuario y la recepción de
mensajes.

La documentación fue una parte importante de este aprendizaje.
Utilicé la guía de Beej como referencia para la comunicación mediante
sockets, las páginas del manual de Linux para comprender las funciones
de memoria y la documentación de Microsoft para desarrollar el cliente.
También incorporé cJSON para interpretar y construir mensajes.
Estos recursos me facilitaron el avance, aunque adaptar sus ejemplos
me obligó a comprender cómo utilizarlos y cómo relacionarlos con el
código que ya tenía
\cite{beej,malloc-man,tcp,json,cjson}.

Al principio, mi prioridad fue conseguir que cada nueva función
trabajara. Esto me permitió avanzar poco a poco, pero también dejó
en evidencia que no había planeado bien las responsabilidades.
Varias tareas terminaron concentradas en \archivo{servidor.c}, y
agregar cambios se volvió cada vez más difícil. Cuando reconocí
este problema, revisé el diseño y procuré aplicar mejor la idea de
``divide y vencerás'', separando tareas y reutilizando código que
había comenzado a repetir.

La versión final comprende las operaciones básicas del chat y un
cliente de consola. Prioricé completar esa interacción y conservé
el servidor monohilo ante la cercanía de la entrega. Mi intención
era contar primero con una versión funcional y después incorporar
los cambios más grandes, pero calculé mal el tiempo que necesitaba.
Esta experiencia me permitió identificar qué habría sido conveniente
prever desde el inicio.

\section{Análisis y diseño}

\subsection{Del planteamiento inicial a un servidor con demasiadas responsabilidades}

Al principio separé el proyecto en servidor, cliente y modelos del
chat. Había creado \archivo{modelos_chat.c} y \archivo{modelos_chat.h}
para representar usuarios y salas, pero esos archivos quedaron vacíos
mientras avanzaba con las conexiones. La información de cada usuario
terminó guardándose dentro de \archivo{Cliente}, y más adelante
renombré esos archivos a \archivo{salas.c} y \archivo{salas.h} para
darles una responsabilidad más clara.

Comencé por un programa mínimo, la lectura del puerto y la apertura
del socket. Después añadí la aceptación de conexiones, la recepción
de bytes y la atención de varios clientes. Esta secuencia me permitió
observar resultados pronto y comprobar cada avance.

Sin embargo, no apliqué adecuadamente el principio de ``divide y
vencerás''. Conforme incorporé el manejo de JSON, la identificación,
las consultas y las respuestas, \archivo{servidor.c} acumuló tareas
distintas y superó las 800 líneas. La cantidad de líneas fue una señal
de que algo no iba bien. En ese punto me costaba encontrar dónde
hacer cada cambio y decidí separar las responsabilidades.

Este diagrama muestra cómo estaba organizado el proyecto en esa
etapa, antes de separar las responsabilidades que había acumulado
el servidor. El cliente en C\# todavía no existía, así que utilicé
netcat para probar las conexiones.

\begin{figure}[htbp]
\centering
\begin{tikzpicture}
    \node[modulo] (main) at (0,0) {
        \textbf{main.c}\\
        Argumentos y arranque
    };

    \node[
        modulo,
        text width=8cm,
        minimum height=1.9cm
    ] (srv) at (0,-2.2) {
        \textbf{servidor.c}\\
        Escucha y conexiones\\
        Recepción y mensajes completos\\
        Solicitudes, usuarios y respuestas
    };

    \node[
        modulo,
        text width=4.6cm
    ] (cli) at (-3,-4.7) {
        \textbf{cliente.c / Cliente}\\
        Socket y búferes por conexión
    };

    \node[
        modulo,
        text width=4.6cm
    ] (nc) at (3,-4.7) {
        \textbf{netcat}\\
        Pruebas de conexión
    };

    \draw[flecha] (main)--(srv);
    \draw[flecha] (srv)--(cli);
    \draw[flecha,<->] (srv)--node[etiqueta,right]{TCP}(nc);
\end{tikzpicture}
\caption{Organización del proyecto antes de separar el controlador y el protocolo.}
\end{figure}

Resolví el problema de acumulación de responsabilidades moviendo la
interpretación y construcción de JSON a \archivo{protocolo.c}, y
las decisiones sobre las operaciones del chat a
\archivo{controlador.c}. Dejé en \archivo{servidor.c} el manejo de
conexiones y la recepción de mensajes completos. Así pude cambiar
una operación sin mezclarla con el código encargado de aceptar
conexiones y recibir bytes.

También conservé los comentarios sobre las adaptaciones y reuní en
funciones comunes tareas que comenzaban a repetirse, como encolar
respuestas y enviar eventos a varios usuarios.

En este punto Git fue clave, porque me permitió consultar versiones
anteriores y comparar cómo había cambiado el programa. Cuando una
integración dejó de funcionar, pude regresar a una versión estable
y retomar los cambios poco a poco.

\subsection{Organización después de la refactorización}

Este diagrama muestra cómo quedaron organizadas las partes del
proyecto después de la refactorización y cómo colaboran entre sí.
Repartí el código en archivos con responsabilidades concretas.
\archivo{Cliente} reúne la información de cada conexión y
\archivo{Servidor} administra la escucha y la atención de los
clientes. Esta organización me ayudó a identificar dónde realizar
cada cambio y a evitar que las tareas volvieran a concentrarse en
\archivo{servidor.c}.

\begin{figure}[htbp]
\centering
\begin{tikzpicture}
    \node[modulo,text width=4.8cm] (srv) at (-3.3,0) {
        \textbf{«módulo» servidor.c}\\
        Conexiones y recepción
    };

    \node[modulo,text width=4.8cm] (con) at (-3.3,-2) {
        \textbf{«módulo» controlador.c}\\
        Operaciones del chat
    };

    \node[modulo,text width=4.8cm] (pro) at (-3.3,-4) {
        \textbf{«módulo» protocolo.c}\\
        Validación y construcción de JSON
    };

    \node[modulo,text width=4.8cm] (cli) at (-3.3,-6) {
        \textbf{«estructura» Cliente}\\
        Nombre, estado y búferes
    };

    \node[modulo,text width=4.8cm] (cx) at (3.3,0) {
        \textbf{«clase» ConexionCliente}\\
        Envío y recepción por TCP
    };

    \node[modulo,text width=4.8cm] (cc) at (3.3,-2) {
        \textbf{«clase» ControladorCliente}\\
        Comandos y recepción
    };

    \node[modulo,text width=4.8cm] (pc) at (3.3,-4) {
        \textbf{«clase» ProtocoloCliente}\\
        JSON y mensajes legibles
    };

    \draw[flecha] (srv)--(con);
    \draw[flecha] (con)--(pro);
    \draw[flecha] (con.west)--++(-.6,0)|-(cli.west);
    \draw[flecha] (srv.west)--++(-2,0)|-(cli.west);
    \draw[flecha] (pro)--node[etiqueta,right]{valida estados}(cli);
    \draw[flecha] (cc)--(cx);
    \draw[flecha] (cc)--(pc);
    \draw[flecha,<->] (srv)--node[etiqueta,above]{TCP / JSON}(cx);
\end{tikzpicture}
\caption{Organización de los módulos del servidor y las clases del cliente después de la refactorización.}
\end{figure}

Tomé como referencia la separación de responsabilidades de
Modelo--Vista--Controlador. Los datos de las conexiones y los
usuarios forman parte del modelo, el controlador decide qué
operación ejecutar y la consola permite interactuar con el programa.
Aun así, no logré una separación completa. Por ejemplo,
\archivo{ProtocoloCliente} también prepara el texto que se muestra
en pantalla, y \archivo{protocolo.c} consulta la función de estados
de \archivo{cliente.c}. Son partes que podría organizar mejor.

Estas fueron algunas de las decisiones que tomé durante el desarrollo.

\begin{center}
\small
\begin{tabularx}{\textwidth}{@{}p{3.7cm}X@{}}
    \toprule
    Decisión & Motivo \\
    \midrule

    Cliente de consola &
    Me permitió concentrarme en la comunicación y en los comandos
    antes de trabajar en una interfaz gráfica. \\

    Bibliotecas JSON &
    Utilicé cJSON y System.Text.Json para interpretar y construir
    mensajes, y dejé en mi código las comprobaciones del protocolo. \\

    Arreglo por descriptor &
    Me permitió localizar directamente al cliente asociado con un
    socket. La búsqueda por nombre todavía requiere recorrer
    los usuarios. \\

    Búfer ampliable &
    Lo necesité para conservar fragmentos cuando un mensaje no
    llegaba completo en una sola lectura. \\

    \bottomrule
\end{tabularx}
\end{center}

\subsection{Tarjetas CRC como revisión del diseño}

Preparé estas tarjetas para revisar las responsabilidades de cada
parte del programa y las otras partes con las que necesita colaborar. No las utilicé al comenzar el proyecto, pero al elaborarlas
pude reconocer con más claridad la separación que buscaba durante
la refactorización. Para el código en C las adapté a módulos y
estructuras, y para C\# utilicé las clases del cliente.

\noindent
\crc{Servidor / servidor.c}
    {Abrir la escucha, aceptar conexiones, recibir mensajes y
    cerrar conexiones.}
    {Cliente y controlador.c.}
\hfill
\crc{Cliente / cliente.c}
    {Guardar nombre y estado, administrar búferes y conservar
    los mensajes pendientes de envío.}
    {servidor.c y controlador.c.}

\vspace{7pt}
\noindent
\crc{Protocolo / protocolo.c}
    {Interpretar y validar JSON, consultar campos y construir
    respuestas y eventos.}
    {cJSON, controlador.c y cliente.c para los estados.}
\hfill
\crc{Controlador / controlador.c}
    {Identificar usuarios, buscar destinatarios, ejecutar
    solicitudes y enviar avisos.}
    {Cliente, protocolo.c y servidor.c.}

\vspace{7pt}
\noindent
\crc{ConexionCliente}
    {Establecer la conexión, enviar y recibir mensajes, y
    liberar los recursos de conexión.}
    {TcpClient, StreamReader, Program y ControladorCliente.}
\hfill
\crc{ProtocoloCliente}
    {Construir solicitudes JSON e interpretar las respuestas
    para mostrarlas al usuario.}
    {System.Text.Json, Program y ControladorCliente.}

\vspace{7pt}
\noindent
\crc{ControladorCliente}
    {Leer comandos, coordinar la recepción, mostrar mensajes
    y gestionar la salida.}
    {ConexionCliente, ProtocoloCliente y Console.}
\hfill
\crc{Salas / salas.c (planeado)}
    {Guardar salas, integrantes e invitaciones, retirar
    usuarios y eliminar salas vacías.}
    {Controlador y Cliente.}

\subsection{Dificultades con los mensajes y la memoria}

Seguí la indicación del protocolo del libro de terminar cada
mensaje con un salto de línea, \archivo{\n}, e ignorar las líneas
vacías. Al implementarlo tuve que comprender que una
lectura del socket no siempre contiene un mensaje completo.
Puede traer solo una parte o varios mensajes juntos, así que
necesité conservar los fragmentos hasta encontrar el salto de línea.

Elegí un búfer de 4096 bytes para recibir datos y fijé el límite
del mensaje completo en 1\,048\,576 bytes, conforme a las indicaciones. También tuve que distinguir el salto de línea
del carácter \archivo{\0}, que termina las cadenas en C y no
debe enviarse como parte del mensaje.

Este diagrama muestra el recorrido de un mensaje público y permite relacionar esa dificultad con las responsabilidades que
separé durante la refactorización. Primero reúno el
mensaje completo, después procesarlo y finalmente enviarlo a
los demás clientes.

\begin{figure}[h]
\centering
\begin{tikzpicture}[x=1cm,y=1cm]
\foreach \x/\nombre in {
    0/Cliente emisor,
    4/Servidor,
    8/Controlador,
    12/Cliente receptor
}{
    \node[
        font=\small\bfseries,
        align=center
    ] at (\x,0) {\nombre};

    \draw[dashed,gray] (\x,-.3)--(\x,-3.4);
}

\draw[flecha]
    (0,-.7)--node[etiqueta,above]
    {JSON + salto de línea}(4,-.7);

\draw[flecha]
    (4,-1.4)--node[etiqueta,above]
    {mensaje completo}(8,-1.4);

\draw[flecha]
    (8,-2.1)--node[etiqueta,above]
    {encola evento}(4,-2.1);

\draw[flecha]
    (4,-2.9)--node[etiqueta,above]
    {envía cuando hay escritura disponible}(12,-2.9);
\end{tikzpicture}
\caption{Secuencia simplificada de un mensaje público. Validación JSON y cola del destinatario se resumen dentro de las acciones del controlador; recibir no equivale a enviar inmediatamente.}
\end{figure}

Otra dificultad fue asumir que enviar un mensaje significaba
transmitirlo completo de una sola vez. Para resolverlo guardé
los bytes pendientes y utilicé \archivo{select()} para saber
cuándo podía continuar el envío. Esto amplió el ejemplo inicial
de Beej, que había tomado como base para atender las conexiones.

Al probar el protocolo también aprendí a distinguir los mensajes
que envía el cliente de los que responde el servidor. Por ejemplo,
envié \archivo{DISCONNECTED} cuando correspondía
\archivo{DISCONNECT}, y recibí una respuesta de mensaje inválido.
Hacer estas pruebas me ayudó a comprobar que no bastaba con que
el JSON estuviera bien escrito, también debía representar una
solicitud válida.

El manejo de memoria requirió especial atención. Consulté
\archivo{man 3 malloc} para comprender la reserva, ampliación y
liberación de memoria. Al utilizar \archivo{realloc()}, guardé
primero el resultado en un apuntador temporal para no perder la
reserva anterior si la ampliación fallaba \cite{malloc-man}.

También tuve que revisar qué datos podía conservar después de
liberar un objeto JSON. El nombre del cliente necesitaba una
copia propia, mientras que las respuestas podían liberarse
después de copiarlas a la cola de salida.

Por último, encontré que cerrar una conexión inmediatamente
después de detectar un error podía impedir que el cliente
recibiera la respuesta. Para resolverlo añadí un cierre pendiente,
de modo que el servidor terminara de enviar lo acumulado antes
de cerrar la conexión.

\section{Implementación}

\subsection{Tecnologías y decisiones prácticas}

Implementé el servidor con C11, sockets POSIX y cJSON. De la guía
de Beej adapté ejemplos para aceptar conexiones, atender varios
clientes y ampliar arreglos. A partir de ellos añadí la información
de cada conexión, los límites de los búferes y el procesamiento
del protocolo \cite{beej}.

Incorporé los archivos y la licencia de cJSON al proyecto y adapté
el ejemplo de lectura de campos de su documentación a los mensajes
del chat. Esto me permitió concentrarme en validar los campos y
ejecutar las operaciones, aunque tuve que aprender a liberar los
objetos que construía la biblioteca \cite{cjson}.

Para el cliente utilicé C\# y un proyecto de consola dirigido a
\archivo{net10.0} y me apoyé en la documentación de Microsoft para utilizar
\archivo{TcpClient}, \archivo{StreamReader} y
\archivo{System.Text.Json} \cite{tcp,json}.

Una de las decisiones del cliente fue permitir que el usuario
escribiera texto normal y comandos en español. Por ejemplo,
\archivo{/estado ausente} se convierte en el valor
\archivo{AWAY} del protocolo. Esto surgió para que el usuario no tuviera que escribir JSON para utilizar el chat.

Antes de contar con el cliente, utilicé netcat para verificar
las conexiones y enviar solicitudes manualmente. Así pude
avanzar con el servidor sin depender de tener terminados ambos
programas al mismo tiempo.

Comencé con una configuración de CMake y después cambié a Make,
debido a la complejidad de CMake. Make me resultó
más fácil de manejar, y terminé reuniendo en el mismo Makefile
la compilación y ejecución del servidor y del cliente. Poder
elegir el puerto y la dirección desde esos comandos me facilitó
las pruebas.

\section{Conclusiones}

Logré comunicar programas escritos en dos lenguajes que no había
utilizado antes. Pasar de observar bytes en una terminal a
intercambiar mensajes desde el cliente de consola me permitió
ver el resultado de cada etapa. También aprendí que recibir
datos, reconocer un mensaje y ejecutar una operación son
problemas diferentes, y que separarlos desde el principio
habría simplificado el desarrollo.

\subsection{Dificultades y correcciones durante el desarrollo}

La primera dificultad fue comprender el manejo de memoria y
apuntadores en C. Los descriptores de socket, las reservas
dinámicas y las cabeceras requirieron una atención a la que
no estaba acostumbrado trabajando con Java. Uno de mis errores
fue utilizar \archivo{memcmp()} cuando necesitaba copiar datos
con \archivo{memcpy()}. En otra ocasión omití el asterisco de
un parámetro en una cabecera, por lo que la declaración dejó
de coincidir con la implementación.

También confundí comparación y asignación, olvidé inicializar
algunos atributos y escribí retornos incorrectos. Un mensaje
válido llegó a provocar una desconexión porque había manejado
mal el resultado de una función. Estos errores me hicieron
ver que conseguir que el programa compilara era solo una parte
del trabajo.

Otro error fue intentar acceder a \archivo{cliente->salida}
después de ejecutar \archivo{free(cliente)}. Lo corregí
liberando primero los recursos internos y después la estructura.
Este caso me ayudó a comprender que debía prestar atención
tanto a qué memoria liberaba como al orden en que lo hacía.

Al integrar cambios dejé llaves mal colocadas, funciones
dentro de otras por accidente, código después de
\archivo{return} y definiciones duplicadas. Durante la
refactorización también aparecieron funciones declaradas
pero sin implementación y llamadas que conservaban nombres
anteriores. A partir de estos problemas procuré compilar y
comprobar cada cambio pequeño antes de continuar.

En C\# tuve dificultades de configuración. Creé el proyecto
inicialmente fuera del repositorio, una clase no se compilaba
porque el archivo tenía una extensión incorrecta y olvidé
importar \archivo{System.Text} para utilizar
\archivo{StringBuilder}. Aunque eran errores pequeños,
me hicieron perder tiempo y aprendí a revisar también
los archivos y la configuración antes de cambiar la lógica.

\subsection{Cosas a mejorar}

Resolver la acumulación de responsabilidades fue una de las
mejoras más importantes del proyecto. Separar el controlador
y el protocolo me permitió reunir código repetido y encontrar
con más facilidad dónde hacer los cambios. Sin embargo,
habría sido mejor definir esas responsabilidades antes de
que el servidor creciera tanto. Utilizar las tarjetas CRC
desde el inicio me habría ayudado a revisar esa organización.

Git me permitió conservar versiones funcionales y retomar
el trabajo después de una integración fallida. Mantendría
esa forma de avanzar en una siguiente versión, con cambios
pequeños y comprobaciones frecuentes, pero dedicaría más
tiempo al análisis antes de comenzar a programar.

Mi mala planeación también afectó el alcance de la entrega.
Dejé pendientes las salas completas, la migración a multihilos
y el cierre ordenado del servidor. Quiero retomarlos después
para terminar el proyecto, pero con una revisión previa de
las responsabilidades y del tiempo que requiere cada parte.

\begingroup
\small

\begin{thebibliography}{9}

\bibitem{beej}
Brian ``Beej'' Hall.
\textit{Beej's Guide to Network Programming}.
Guía y ejemplos \archivo{selectserver.c} y \archivo{pollserver.c}.
\url{https://beej.us/guide/bgnet/}.

\bibitem{chat-server}
Yorick de Wid.
\textit{Chat-Server}.
Repositorio de código fuente en GitHub.
\url{https://github.com/yorickdewid/Chat-Server}.

\bibitem{malloc-man}
Linux man-pages project.

\bibitem{cjson}
Dave Gamble y colaboradores.
\textit{cJSON}.
Código, README y licencia.
\url{https://github.com/DaveGamble/cJSON}.

\bibitem{tcp}
Microsoft.
\textit{Uso de TcpClient y TcpListener}.
Microsoft Learn.
\url{https://learn.microsoft.com/es-es/dotnet/fundamentals/networking/sockets/tcp-classes}.

\bibitem{json}
Microsoft.
\textit{Uso de un DOM JSON en System.Text.Json}.
Microsoft Learn.
\url{https://learn.microsoft.com/es-es/dotnet/standard/serialization/system-text-json/use-dom}.

\end{thebibliography}

\endgroup
\end{document}
